\documentclass{article}
\usepackage[utf8]{inputenc}
\usepackage{amsmath, amssymb}
\usepackage[T1]{fontenc}
\title{Mathematik Dokument}
\author{Anh Le Xuan, Max Berthold, Moritz Engelmann}
\date{\today}

\begin{document}
\maketitle

\section{Einleitung}
Das ist die mathematische Dokumention für den LogisticsLab Kurs WiSe2025.

\section{euklidische Distanz}
Wir haben eine Menge von Maschinen, jede mit x- und y-Koordinaten.
Wobei M die Menge aller Maschinen, $m_i$ die Maschinen ID, $x_i$ die x-Koordinate und $y_i$ die y-Koordinate der Maschine $i$ ist:
\begin{equation}
  M = \{M_i\ | \ i = 1, 2, \ldots, n\}
\end{equation}
\begin{equation}
  M = \{(m_i, x_i, y_i) \ | \ i = 1, 2, \ldots, n\}
\end{equation}
In unserem Fall sieht die Menge der Maschinen $M$ laut der Textdatei "machine\_positions.txt" folgendermaßen aus:
\begin{equation}
  M = \{M_1, M_2, M_3, \ldots\}
\end{equation}
\begin{equation}
  M = \{(1, 30, 5), (2, 30, 35), (3, 60, 35), \ldots\}
\end{equation}
Die euklidische Distanz $d: M \times M \rightarrow \mathbb{R} $ zwischen zwei Maschinen $M_i$ und $M_j$ wird berechnet durch:
\begin{equation}
  d_{ij}= d(m_i, m_j) = \sqrt{(x_i - x_j)^2 + (y_i - y_j)^2}
\end{equation}
Allgemein als Hinweis gilt, dass das Kreuzprodukt kartesisches Produkt zweier Mengen $A$ und $B$ definiert ist als:
\begin{equation}
  A \times B = \{(a, b) \ | \ a \in A, b \in B\}
\end{equation}
Beispiel: $A = \{1, 2\}, B = \{3, 4\} \implies A \times B = \{(1, 3), (1, 4), (2, 3), (2, 4)\}$ \\ \\ \\ \\ \\
In unserem konkreten Fall mit 3 Maschinen: \\
$M_1 = (1,1)$ \\
$M_1 = (4,1)$ \\
$M_1 = (4,5)$ \\
$M \times M = \{(M_1, M_1), (M_1, M_2), (M_1, M_3), (M_2, M_1), (M_2, M_2), (M_2, M_3), (M_3, M_1), (M_3, M_2), (M_3, M_3)\}$
Distanzberechnung: \\
$d(M1,M1) = 0$ \\
$d(M1,M2) = \sqrt{(1-4)^2 + (1-1)^2} = 3$ \\
$d(M1,M3) = \sqrt{(1-4)^2 + (1-5)^2} = 5$ \\
usw.\\
Darstellung als Matrix:
\begin{equation}
  D = \begin{bmatrix}
    d(M_1, M_1) & d(M_1, M_2) & d(M_1, M_3) \\
    d(M_2, M_1) & d(M_2, M_2) & d(M_2, M_3) \\
    d(M_3, M_1) & d(M_3, M_2) & d(M_3, M_3)
  \end{bmatrix}
  =
  \begin{bmatrix}
    0 & 3 & 5 \\
    3 & 0 & 4 \\
    5 & 4 & 0
  \end{bmatrix}
\end{equation}








\end{document}